\section{Building the code and setting up a case}
\label{app:usage}

\subsection{Building}

\NHFOAM\ needs the ESI OpenFOAM line (\url{https://www.openfoam.com}). It was
compiled and run with v1912, and every API call was checked against the v2512
headers. With the OpenFOAM environment loaded, \code{./Allwmake} builds
\begin{itemize}[nosep]
  \item \code{libNHFOAM}: the $\sigma$-column structure, mesh motion, the
        interface pressure matrix and the \code{sigmaColumn} preconditioner,
        linear wave theory, relaxation zones, seabed motion and the
        \code{nhMovingSlipVelocity} condition;
  \item \code{NHFOAM}: the solver;
  \item \code{nhSetWaves}: initial conditions (\code{standing},
        \code{solitary}, \code{linear}, \code{circularMode}).
\end{itemize}

\subsection{Mesh and fields}

The mesh must be a vertical extrusion between a top patch (\code{surface}) and
a bottom patch (\code{bottom}), with the points of each column on a vertical
line: \code{blockMesh} with vertical block edges, \code{extrudeMesh} from a 2D
mesh, or gmsh (\code{Extrude} with \code{Layers} and \code{Recombine}) followed
by \code{gmshToFoam}. The bathymetry is built into the mesh, with the top of
the mesh at the still-water level. Uniform layers are best with the interface
pressure (Section~\ref{sec:dispersion}): two layers keep the period error within
0.6\% up to $kh = 7$, three layers within 0.21\%.

\begin{table}[ht]
\centering
\caption{Boundary conditions of the fields in \code{0/}.}
\label{tab:fields}
\small
\begin{tabular}{llll}
\toprule
field & \code{surface} & \code{bottom}, walls & notes \\
\midrule
\code{U} & \code{zeroGradient} & \code{slip} or \code{noSlip} & \code{nhMovingSlipVelocity} on a moving bottom \\
\code{p\_rgh} & \code{mixed} & \code{fixedFluxPressure} & the mixed coefficients are set by \NHFOAM \\
\code{eta} & \code{zeroGradient} & \code{zeroGradient} & one value per column \\
\bottomrule
\end{tabular}
\end{table}

\subsection{\code{constant/NHFOAMProperties}}

\begin{verbatim}
surfacePatch     surface;     // top patch
bottomPatch      bottom;      // bottom patch
minDepth         1e-4;        // smallest depth used to place mesh points
stillWaterLevel  0;           // reference for the relaxation-zone targets
pressureLocation interface;   // or cell

gauges                        // -> postProcessing/gauges/<t0>/eta.dat
{
    G1  (5.7 0 0);
}

forcingZones
{
    inlet                     // generation and absorption
    {
        type        wave;
        start       (8 0 0);  // sigma_r = 0 here ...
        end         (0 0 0);  // ... sigmaMax here
        sigmaMax    20;  exponent 2;
        waveHeight  0.02;  wavePeriod 2.02;  depth 0.4;
        direction   (1 0 0);  rampTime 4.04;
    }
    outlet
    {
        type        sponge;
        start       (46 0 0);  end (54 0 0);
        sigmaMax    20;
        // relaxEta false;    // damp the velocity only
    }
}

breaking { model HFA; alpha 0.6; beta 0.3; hydrostaticFactor 1000; }

bottomMotion                  // uplift, or slide (see the source)
{
    type uplift;  min (-0.305 -1 0);  max (0.305 1 0);
    amplitude 0.0005;  history exponential;  riseTime 0.1;
}
\end{verbatim}

\subsection{Schemes and solution controls}

\begin{verbatim}
ddtSchemes { default backward 1; }        // "1": full ddt flux correction
divSchemes
{
    div(phi,U)  Gauss filteredLinear2V 0.2 0;  // non-breaking waves
//  div(phi,U)  Gauss limitedLinearV 1;        // with breaking (bores)
}

PIMPLE
{
    nOuterCorrectors          1;
    nCorrectors               3;
    nNonOrthogonalCorrectors  1;     // over sloping beds
    fluxExtrapolation         true;  // default with backward
}

p_rgh
{
    solver          PCG;
    preconditioner  sigmaColumn;     // DIC in the cell layout
    tolerance       1e-10;
    relTol          0.01;
}
\end{verbatim}

Without the ddt coefficient OpenFOAM limits its ddt flux correction, and waves
are visibly damped; \NHFOAM\ warns if it is missing. In the interface layout
only PCG and PBiCGStab are used, and any preconditioner other than
\code{diagonal} or \code{none} is replaced by \code{sigmaColumn}. The unknowns
are written as \code{pI}, the pressure on the bottom interface of each cell.

The time step should resolve the shortest wave that matters, not only the
peak period: behind the submerged bar the free higher harmonics have periods
of $T/2$ and $T/3$ (Section~\ref{sec:bar10}).

\subsection{Initial conditions and parallel runs}

Initial conditions are set with \code{nhSetWaves}
(\code{system/nhSetWavesDict}); \NHFOAM\ moves the mesh onto the initial
$\eta$ at start-up. Whole columns must stay on one processor, so the
decomposition uses the \code{structured} method on the bottom patch:
\begin{verbatim}
method structured;
structuredCoeffs
{
    method  simple;
    patches (bottom);
    simpleCoeffs { n (2 1 1); }
}
\end{verbatim}

\subsection{Tutorials and validation scripts}

The tutorials are \code{standingWave}, \code{solitaryWave},
\code{submergedBar}, \code{circularBasin} (3D unstructured prisms),
\code{bottomUplift}, \code{landslide3D} and \code{breakingSlope}. Each runs with
\code{./Allrun}. The folder \code{validation/} holds the scripts that build the
validation cases (\code{make\_case.py}, \code{ma2012.py},
\code{other\_models.py} for the NHWAVE and REEF3D::NHFLOW input) and make the
figures of this report (\code{make\_figures.sh}).
